原子核（9.14\(\sim\)9.20）

\lsn{1}填写下面物理史内容：

\ding{172}\lunl{4em}提出了原子论的基本假设，认为物质是由微小的、不可分割的粒子组成的，这些粒子称为原子。

\ding{173}提出枣糕模型的物理学家是\underline{\hbox to 5em{}}，他发现了\lunl{3em}

\ding{174}\underline{\hbox to 4em{}}通过\(\alpha\)粒子散射实验提出了原子的\underline{\hbox to 5em{}}模型

\ding{175}\underline{\hbox to 5em{}}最先发现了天然放射现象

\ding{176}最先发现质子的物理学家是\underline{\hbox to 4em{}}，发现中子的物理学家是\underline{\hbox to 5em{}}

\ding{177}发现\(X\)射线的是\underline{\hbox to 3em{}}

\ding{178}第一次在实验室实现人工转变的物理学家是\underline{\hbox to 4em{}}

\lsn{2}填写关于三种放射线的内容：

\begin{table}[htbp]
	\begin{tblr}{c|p{5em}|p{5em}|p{5em}}
	\hline
	&\(\alpha\) 射线&\(\beta\)射线&\(\gamma\)射线\\
	\hline
	组成&&&\\
	\hline
    电荷量&&&\\
    \hline
	速度&&&\\
	\hline
	贯穿本领&&&\\
	\hline
	电离能力&&&\\
	\hline
	\end{tblr}
\end{table}

\lsn{3}已知\,\qty{1}{u}=\qty{1.6605402e-27}{kg}，\(c =\) \,\qty{3e8}{m/s}，\(e =\) \,\qty{1.6e-19}{C}，请计算：\,\qty{1}{u}的质量能量等价值\(E\)是多少\,\qty{}{J}？多少\,\qty{}{MeV}?并写出计算过程。
\vspace*{3cm}

\lsn{4}写出常见的粒子表示符号：

质子：\lunl{2em}中子:\lunl{2em}氢的三种同位素：\lunl{2em}、\lunl{2em}、\lunl{2em}

氦的两种同位素：\lunl{2em}、\lunl{2em}碳的两种同位素：\lunl{2em}、\lunl{2em}

正负电子\lunl{2em}、\lunl{2em}中微子\lunl{2em}

\lsn{5}写出\(\alpha\)衰变一般性方程与本质方程：
\vspace*{2cm}

\lsn{6}写出\(\beta\)衰变一般性方程与本质方程：

\newpage
\lsn{7}关于半衰期有哪些考点？
\vspace*{3cm}

\lsn{8}\lkf{2022全国甲}
	两种放射性元素的半衰期分别为\(t_0\)和\(2t_0\)，\(t = 0\)时刻这两种元素的原子核总数为\(N\)，在\(t=2t_0\)时刻，尚未衰变的原子核总数为\(\frac{N}{3}\)，则在\(t=4t_0\)时刻，尚未衰变的原子核总为\lunl{3em}

\lsn{9}\lkf{2023福建，7}
	福建福清核电站采用我国完全自主研发的“华龙一号”反应堆技术，建设了安全级别世界最高的机组。机组利用\ce{^{235}U}核裂变释放的能量发电，典型的核反应方程为\ce{^1_0n + ^{235}_{92}U \to ^{141}_{Z}Ba + ^{A}_{36}Kr + 3^1_0n}，则\(A =\) \underline{\hbox to 2em{}};\(Z =\) \underline{\hbox to 2em{}}，若核反应过程中质量亏损\,\qty{1}{\g}，释放的能量为\underline{\hbox to 3em{}}J。(光速大小取\,\qty{3.0e8}{m/s})

\lsn{10}\lkf{2024全国甲，14，6分}氘核可通过一系列聚变反应释放能量，总的反应效果可用\ce{6^2_1H\to 2 ^4_2He + x^1_0n +y^1_1p +43.15MeV}表示，式中\(x =\)\lunl{2em}、\(y =\)\lunl{2em}

\lsn{11}写出下面两个典型的人工核反应方程：

\(\alpha\)粒子轰击\ce{^{14}N}发现质子：\lunl{20em}

\(\alpha\)粒子轰击\ce{^{9}Be}发现中子：\lunl{20em}

\lsn{12}所有的核反应方程都必须满足哪些守恒定律？
\vspace*{2cm}

\lsn{13}怎样理解结合能与比结合能这两个概念？
\vspace*{6cm}

\lsn{14}阐述比结合能与原子核稳定性的关系。